\documentclass[11pt]{article}
\usepackage{amsmath, amssymb}
\usepackage{booktabs}
\usepackage{geometry}
\usepackage{setspace}
\usepackage{hyperref}

\geometry{margin=1in}
\onehalfspacing

\begin{document}

\title{Replication Summary: Szymanowska et al.\ (2014)}
\author{Garbuzov, Mitchell}
\date{\today}
\maketitle

\section{Overview}

This document summarizes the replication of the main summary statistics for
commodity futures returns in Szymanowska et al.\ (2014). The focus is on
Table~1--style results: annualized mean returns, standard deviations, and
Newey--West $t$-statistics for Short Roll (SR) and Excess Holding (EH) returns
across commodity sectors.

% NOTE TO SELF:
% If you later replicate Table 2 or other tables, you can:
%   - generate separate .tex table files, and
%   - include them with additional \input{...} commands and sections.

\section{Data and Sample}

The replication uses bimonthly futures data for 21 commodities, grouped into
seven sectors (Energy, Meats, Metals, Grains, Oilseeds, Softs, and industrial
materials), over the period March 1986 to December 2010. For each observation
date, prices of up to four nearby maturities are used to construct log returns
for Short Roll and Excess Holding strategies.

\subsection{Short Roll and Excess Holding Returns}

Let $F_t^{(n)}$ denote the futures price with maturity $n$ (in months) at time
$t$, and let the bimonthly observation interval be $\Delta t$ (six periods per
year).

\paragraph{Short Roll (SR).}
The Short Roll strategy continually rolls into contracts with fixed
time-to-maturity $n$, rather than holding a single contract to maturity. The
log Short Roll return from $t$ to $t+\Delta t$ for maturity $n$ is, at a high
level,
\[
  r^{\text{SR},n}_{t\to t+\Delta t}
    = \ln F_{t+\Delta t}^{(n)} - \ln F_t^{(n)}.
\]
In the implementation, these returns are computed at the commodity level and
labelled $\texttt{sr\_1}, \dots, \texttt{sr\_4}$; sector-level SR indices are
formed as equal-weighted averages across commodities in each sector.

\paragraph{Excess Holding (EH).}
Excess Holding returns compare a buy-and-hold strategy over $n$ bimonthly
periods to the Short Roll strategy over the same horizon. Let
$r^{\text{Hold},n}_{t\to t+n}$ be the log return from buying an $n$-maturity
contract at $t$ and holding it until $t+n$. Then the Excess Holding return is
\[
  r^{\text{EH},n}_{t\to t+n}
    = r^{\text{Hold},n}_{t\to t+n} - r^{\text{SR},n}_{t\to t+n}.
\]
These are stored as $\texttt{eh\_2}, \texttt{eh\_3}, \texttt{eh\_4}$ and
aggregated to the sector level via equal-weighted averages.

\subsection{Annualization and Newey--West $t$-statistics}

All summary statistics are based on bimonthly log returns. There are six
bimonthly periods per year. For a strategy with holding period
$n_{\text{periods}}$ (in bimonthly units), the annualization factor is
\[
  \text{ann\_factor} = \frac{6}{n_{\text{periods}}}.
\]
Given the sample mean $\mu_{\text{raw}}$ and standard deviation
$\sigma_{\text{raw}}$ of the bimonthly returns, the annualized mean and
standard deviation are
\[
  \mu_{\text{ann}} = \mu_{\text{raw}} \times \text{ann\_factor}, \qquad
  \sigma_{\text{ann}} = \sigma_{\text{raw}} \times \sqrt{\text{ann\_factor}}.
\]
Short Roll returns use $n_{\text{periods}} = 1$ for all maturities
$n = 1,2,3,4$, while Excess Holding returns use
$n_{\text{periods}} = n$ (e.g.\ $n_{\text{periods}} = 2$ for $\texttt{eh\_2}$).

To account for autocorrelation and heteroskedasticity in the bimonthly series,
$t$-statistics are computed using a Newey--West HAC estimator. For each return
series $\{r_t\}_{t=1}^T$, we estimate
\[
  r_t = \alpha + \varepsilon_t,
\]
and report the Newey--West $t$-statistic
\[
  t_{\text{NW}} = \frac{\hat{\alpha}}{\text{SE}_{\text{NW}}(\hat{\alpha})},
\]
with a lag length chosen as the maximum of
$n_{\text{periods}} - 1$ and a rule-of-thumb $\lfloor T^{1/3} \rfloor$.

\section{Summary Statistics for SR and EH}

Table~\ref{tab:table1_sr_eh} reports the annualized mean, standard deviation,
and Newey--West $t$-statistics for Short Roll (Panel A) and Excess Holding
(Panel B) returns across sectors, as well as for an equally-weighted index
across all commodities.

\begin{table}[htbp]
  \centering
  \caption{Summary Statistics for Short Roll and Excess Holding Returns}
  \label{tab:table1_sr_eh}
  \footnotesize
  %
  % NOTE:
  % Generate a table file (e.g. table1_sr_eh.tex) from Python and place it in
  % the same directory as this .tex file, then input it here.
  %
  \begin{table}
\caption{Summary statistics for Short Roll and Excess Holding returns}
\label{tab:table1_sr_eh}
\begin{tabular}{lrrrrrrrrrrrrl}
\toprule
sector & mean_ann_n1 & std_ann_n1 & t_stat_n1 & mean_ann_n2 & std_ann_n2 & t_stat_n2 & mean_ann_n3 & std_ann_n3 & t_stat_n3 & mean_ann_n4 & std_ann_n4 & t_stat_n4 & panel \\
\midrule
Energy & 0.0888 & 0.3461 & 1.3206 & 0.0886 & 0.3305 & 1.3691 & 0.0882 & 0.3162 & 1.4086 & NaN & NaN & NaN & Short Roll \\
Meats & 0.0039 & 0.1296 & 0.1300 & 0.0365 & 0.1171 & 1.3294 & 0.0281 & 0.1101 & 1.1878 & NaN & NaN & NaN & Short Roll \\
Metals & 0.0464 & 0.1747 & 1.2343 & 0.0415 & 0.1760 & 1.0889 & 0.0362 & 0.1751 & 0.9371 & NaN & NaN & NaN & Short Roll \\
Grains & -0.0716 & 0.1902 & -1.5470 & -0.0235 & 0.1873 & -0.4837 & -0.0394 & 0.1748 & -0.8501 & NaN & NaN & NaN & Short Roll \\
Oilseeds & 0.0124 & 0.2366 & 0.2853 & 0.0267 & 0.2320 & 0.6096 & 0.0402 & 0.2214 & 0.9308 & NaN & NaN & NaN & Short Roll \\
Softs & -0.0782 & 0.1902 & -2.0915 & -0.0715 & 0.1811 & -1.9486 & -0.0700 & 0.1761 & -1.9106 & NaN & NaN & NaN & Short Roll \\
Ind_Materials & -0.0584 & 0.1980 & -1.2757 & -0.0187 & 0.1728 & -0.4789 & -0.0090 & 0.1540 & -0.2610 & NaN & NaN & NaN & Short Roll \\
EW_All & -0.0162 & 0.1168 & -0.5914 & 0.0053 & 0.1148 & 0.1936 & 0.0023 & 0.1106 & 0.0880 & NaN & NaN & NaN & Short Roll \\
Energy & NaN & NaN & NaN & -0.0001 & 0.0230 & -0.0395 & -0.0002 & 0.0366 & -0.0482 & NaN & NaN & NaN & Excess Holding \\
Meats & NaN & NaN & NaN & 0.0163 & 0.0327 & 3.8422 & 0.0195 & 0.0430 & 3.0782 & NaN & NaN & NaN & Excess Holding \\
Metals & NaN & NaN & NaN & -0.0037 & 0.0064 & -4.2160 & -0.0065 & 0.0089 & -4.4414 & NaN & NaN & NaN & Excess Holding \\
Grains & NaN & NaN & NaN & 0.0240 & 0.0291 & 5.0405 & 0.0264 & 0.0387 & 3.9611 & NaN & NaN & NaN & Excess Holding \\
Oilseeds & NaN & NaN & NaN & 0.0071 & 0.0200 & 2.4134 & 0.0141 & 0.0342 & 2.8919 & NaN & NaN & NaN & Excess Holding \\
Softs & NaN & NaN & NaN & 0.0033 & 0.0169 & 1.2089 & 0.0053 & 0.0256 & 1.1990 & NaN & NaN & NaN & Excess Holding \\
Ind_Materials & NaN & NaN & NaN & 0.0199 & 0.0402 & 2.7012 & 0.0306 & 0.0550 & 2.8717 & NaN & NaN & NaN & Excess Holding \\
EW_All & NaN & NaN & NaN & 0.0106 & 0.0122 & 4.8123 & 0.0135 & 0.0180 & 4.0839 & NaN & NaN & NaN & Excess Holding \\
\bottomrule
\end{tabular}
\end{table}

\end{table}

\section{Adding Further Results}

This document is designed to be extensible. To add more content:

\begin{itemize}
  \item \textbf{Additional tables:}
        Generate new LaTeX table files from Python
        (e.g.\ \texttt{table2_basis_sorts.tex}) and include them with
        \verb|\input{table2_basis_sorts.tex}| in new sections.
  \item \textbf{Figures:}
        Save plots from your code (e.g.\ PNG or PDF) into the
        \texttt{reports/} or \texttt{output/} directory and include them with
        \verb|\includegraphics|.
        For example:
\begin{verbatim}
\begin{figure}[htbp]
  \centering
  \includegraphics[width=0.7\textwidth]{basis_factor_plot.pdf}
  \caption{Basis-sorted high-minus-low portfolio returns.}
\end{figure}
\end{verbatim}
  \item \textbf{Discussion sections:}
        Add new \verb|\section| or \verb|\subsection| blocks to interpret
        differences between your replication and the published results, or to
        document robustness checks.
\end{itemize}

% TODO:
% - Add a section on Table 2 (basis sorts) once the replication is complete.
% - Add plots of time series (e.g., EW_All SR_1 over time) as figures.

\end{document}
